\subsection{Perceived AI Support — AI Questionnaire}
\label{sec:eval-ai-support}

The SUS (\S\ref{sec:eval-sus}) and R-TLX (\S\ref{sec:eval-rtlx})
characterise the interface as it was experienced; they say nothing
directly about whether participants felt the AI itself had been
useful. The Perceived AI Support questionnaire targets exactly this
gap. After each interaction mode, in addition to the SUS and R-TLX,
participants completed a six-item Perceived AI Support questionnaire
on a five-point Likert scale (1 = strong disagreement, 5 = strong
agreement). The six items capture distinct facets of how participants
perceived the AI's contribution to the document they had just
created: helpfulness (item~1), efficiency (item~2), transparency
(item~3), trust (item~4), legal-basis confidence (item~5), and
re-use intent (item~6). Item~5 is the primary self-report confidence
measure referenced in Chapter~\ref{ch:methodology}. The aggregated
``All'' column in the tables below pools the three per-mode samples
into a single per-item descriptive across all $207$ ratings. Given
the consistent Form-first ordering on usability and workload, the
question of interest here is whether perceived AI value follows the
same pattern or breaks from it.

The six items are reported in their administered order, beginning
with broad helpfulness and tightening towards specific facets such
as transparency, trust, and re-use intent. Table~\ref{tab:ai-per-item}
reports the mean response per item per mode and the aggregated mean
across the three modes; the full descriptive distribution per item
is reported in Table~\ref{tab:ai-per-item-distribution}.
Figure~\ref{fig:ai-subquestions} shows the corresponding box-plot
distributions.

\begin{table}[htbp]
\centering
\caption{Perceived AI Support, mean response per item by mode
($n = 69$ per mode), five-point Likert scale (higher is better). The
aggregated ``All'' column is the unweighted mean of the three per-mode
means.}
\label{tab:ai-per-item}
\small
\begin{tabularx}{\linewidth}{@{}c X r r r r@{}}
\toprule
\# & Item & Form & Ask & Chat & All \\
\midrule
1 & The AI was helpful in creating this document. & 4.23 & 4.32 & 4.13 & 4.23 \\
2 & The AI helped me to complete the document faster. & 4.43 & 4.25 & 3.99 & 4.22 \\
3 & I understood why the AI suggested specific GDPR fields. & 3.51 & 4.16 & 3.75 & 3.81 \\
4 & I trust that the AI made correct suggestions. & 3.41 & 3.80 & 3.68 & 3.63 \\
5 & I am confident that the AI identified the correct legal basis. & 3.58 & 3.74 & 3.81 & 3.71 \\
6 & I would use this mode again. & 3.99 & 3.55 & 3.42 & 3.65 \\
\bottomrule
\end{tabularx}
\end{table}

\begin{sidewaysfigure}
\centering
\includegraphics[width=\linewidth]{ai_subquestions_boxplot.png}
\caption{Box plots of the six Perceived AI Support items by mode
($n = 69$ per mode). Boxes show the interquartile range; the central
line marks the median, and the diamond marks the mean.}
\label{fig:ai-subquestions}
\end{sidewaysfigure}

Participants find the AI most helpful in Ask, breaking the
Form-first pattern seen on the usability instruments. For
\emph{``The AI was helpful in creating this document''} (item~1),
Form scores $4.23$, Ask $4.32$, Chat $4.13$; aggregated $4.23$.
Helpfulness is rated highest in Ask and lowest in Chat, with a
small absolute spread of roughly $0.20$~points.

When asked specifically about speed, participants give the most
credit to Form and the least to Chat — the felt speed-up shrinks
the more conversational the mode becomes. For \emph{``The AI helped
me to complete the document faster''} (item~2), Form is rated
highest ($4.43$), Ask second ($4.25$), Chat lowest ($3.99$);
aggregated $4.22$. Perceived speed-up decreases monotonically with
increasing LLM involvement, dropping by approximately
$0.45$~points from Form to Chat.

Ask makes participants feel they understand \emph{why} the AI
suggests what it does, by a clear margin — the widest reversal of
the Form-first ordering across the six items. For \emph{``I
understood why the AI suggested specific GDPR fields''} (item~3),
Form scores lowest ($3.51$), Chat second ($3.75$), Ask highest
($4.16$); aggregated $3.81$. Transparency is rated approximately
$0.65$~points higher in Ask than in Form, the only mode that
explicitly explains each suggestion.

Trust in the AI's suggestions tracks transparency: participants
trust the LLM-driven modes more than Form, with Ask leading. For
\emph{``I trust that the AI made correct suggestions''} (item~4),
Form scores lowest ($3.41$), Chat second ($3.68$), Ask highest
($3.80$); aggregated $3.63$. Trust increases by approximately
$0.27$--$0.39$~points from Form to the LLM-driven modes.

Participants feel most confident that the legal basis is correct in
Chat, and least confident in Form — confidence rises monotonically
with LLM involvement. For \emph{``I am confident that the AI
identified the correct legal basis''} (item~5, primary confidence
measure), the three modes show a monotonic increase: Form $3.58$,
Ask $3.74$, Chat $3.81$; aggregated $3.71$. Confidence in legal
correctness rises by approximately $0.23$~points from Form to
Chat.

When asked which mode they would actually use again, participants
return to the Form-first ordering of the SUS — and the gap is the
widest of the six AI-support items. For \emph{``I would use this
mode again''} (item~6), Form scores highest ($3.99$), Ask second
($3.55$), Chat lowest ($3.42$); aggregated $3.65$. Re-use intent
decreases by approximately $0.57$~points from Form to Chat, the
largest negative gap of the six items.

Beyond the per-mode means, the full descriptive distribution shows
how concentrated or spread participants' responses were on each
item. Table~\ref{tab:ai-per-item-distribution} reports the mean,
standard deviation, median, Q1, and Q3 per item per mode. The minimum
is $1$ and the maximum is $5$ for every cell (both endpoints of the
Likert scale are realised in every item~$\times$~mode combination),
so they are omitted from the table for compactness.

\begin{table}[htbp]
\centering
\caption{Perceived AI Support, full per-item descriptive distribution
by mode ($n = 69$ per mode). Source:
\texttt{ai\_support\_per\_item\_summary.csv}.}
\label{tab:ai-per-item-distribution}
\small
\begin{tabular}{@{}c l r r r r r@{}}
\toprule
\# & Mode & Mean & SD & Median & Q1 & Q3 \\
\midrule
\multirow{4}{*}{1} & All  & 4.23 & 0.91 & 4 & 4 & 5 \\
                   & Form & 4.23 & 0.97 & 5 & 4 & 5 \\
                   & Ask  & 4.32 & 0.76 & 4 & 4 & 5 \\
                   & Chat & 4.13 & 1.00 & 4 & 4 & 5 \\
\midrule
\multirow{4}{*}{2} & All  & 4.22 & 1.13 & 5 & 4 & 5 \\
                   & Form & 4.43 & 0.96 & 5 & 4 & 5 \\
                   & Ask  & 4.25 & 1.03 & 5 & 4 & 5 \\
                   & Chat & 3.99 & 1.32 & 4 & 4 & 5 \\
\midrule
\multirow{4}{*}{3} & All  & 3.81 & 1.06 & 4 & 3 & 5 \\
                   & Form & 3.51 & 1.17 & 4 & 3 & 4 \\
                   & Ask  & 4.16 & 0.92 & 4 & 4 & 5 \\
                   & Chat & 3.75 & 0.98 & 4 & 3 & 5 \\
\midrule
\multirow{4}{*}{4} & All  & 3.63 & 0.97 & 4 & 3 & 4 \\
                   & Form & 3.41 & 1.09 & 4 & 3 & 4 \\
                   & Ask  & 3.80 & 0.87 & 4 & 3 & 4 \\
                   & Chat & 3.68 & 0.92 & 4 & 3 & 4 \\
\midrule
\multirow{4}{*}{5} & All  & 3.71 & 0.88 & 4 & 3 & 4 \\
                   & Form & 3.58 & 0.90 & 4 & 3 & 4 \\
                   & Ask  & 3.74 & 0.87 & 4 & 3 & 4 \\
                   & Chat & 3.81 & 0.88 & 4 & 3 & 4 \\
\midrule
\multirow{4}{*}{6} & All  & 3.65 & 1.25 & 4 & 3 & 5 \\
                   & Form & 3.99 & 1.14 & 4 & 4 & 5 \\
                   & Ask  & 3.55 & 1.16 & 4 & 3 & 4 \\
                   & Chat & 3.42 & 1.39 & 4 & 2 & 5 \\
\bottomrule
\end{tabular}
\end{table}

An aggregated Friedman test on the per-mode aggregate is non-significant
($\chi^2(2) = 1.05$, $p = 0.592$, $W = 0.008$), indicating that
perceived AI support does not differ meaningfully across modes when
the six items are considered together. The single legal-basis
confidence item~$5_{ai}$ is also non-significant under Friedman
($\chi^2(2) = 3.76$, $p = 0.152$). The re-use item~$6_{ai}$ does
reach significance ($\chi^2(2) = 11.12$, $p = 0.004$, $W = 0.081$),
with a large rank-biserial effect for Form versus Chat
($r = +0.543$). Full statistics are reproduced in
Appendix~\ref{app:inferential}.

Perceived AI support thus splits the six items into two camps: items
that frame the AI as an explainer (transparency, trust, legal-basis
confidence) favour the LLM-driven modes, while items that frame it
as a productivity aid (efficiency, re-use intent) favour Form. The
behavioural counterpart of these per-item ratings — which mode
participants actually choose when asked to rank — is examined next
in \S\ref{sec:eval-freeform}.
